% Generated draft: requires mathematical and source review.
\documentclass[11pt]{article}
\usepackage[margin=25mm]{geometry}
\usepackage{fontspec}
\setmainfont{Latin Modern Roman}
\usepackage{amsmath,amssymb,mathtools,unicode-math}
\setmathfont{Latin Modern Math}
\usepackage{enumitem,needspace,xurl,hyperref}
\setlength{\parindent}{0pt}
\setlength{\parskip}{5pt}
\setlength{\emergencystretch}{3em}
\newcommand{\N}{\mathbb{N}}
\newcommand{\Z}{\mathbb{Z}}
\newcommand{\Q}{\mathbb{Q}}
\newcommand{\R}{\mathbb{R}}
\newcommand{\C}{\mathbb{C}}
\newcommand{\F}{\mathbb{F}}
\newcommand{\A}{\mathbb{A}}
\newcommand{\PP}{\mathbb{P}}
\newcommand{\E}{\mathbb{E}}
\newcommand{\eps}{\varepsilon}
\newcommand{\dd}{\,\mathrm{d}}
\begin{document}
\section*{UnsolvedMath ID: 30001305}
\textbf{Source problem code:} catalogue entry
\subsection{Definitions and mathematical statement}
Let \(F\) be a totally real number field, and let \(\Sigma_\infty\) denote its real places. Fix a Hilbert newform \(f\) over \(F\). Write \(\mathcal B(f)\) provisionally for the collection of quaternion algebras \(B/F\) over which the source intends to vary and for which the Jacquet--Langlands transfer denoted by \(f_B\) exists. The source does not specify this collection completely, the hypotheses on \(f\), or the choice and normalization of a form representing each transfer.

For \(x,y\in\C\), define
\[
 x\sim_{\overline{\Q}^{\times}}y
 \quad\Longleftrightarrow\quad
 x=a y\text{ for some }a\in\overline{\Q}^{\times},
\]
where \(\overline{\Q}\subset\C\) is the field of algebraic numbers. Let
\[
 S(B)=\{v\in\Sigma_\infty:B\otimes_{F,v}\R\cong M_2(\R)\}
\]
be the set of real places at which \(B\) splits. This interprets the source's condition \(v\nmid\operatorname{disc}(B)\) as being unramified at the real place \(v\). The source does not specify a discriminant convention that includes infinite places.

\emph{Extracted Shimura period question.} With the admissible family, transfer representatives, and pairings fixed in the intended source-compatible manner, do there exist constants \((c_v)_{v\in\Sigma_\infty}\in\C^{\Sigma_\infty}\), depending on \(F\) and \(f\) but independent of \(B\), such that
\[
 \forall B\in\mathcal B(f),\qquad
 \langle f_B,f_B\rangle
 \sim_{\overline{\Q}^{\times}}
 \prod_{v\in S(B)}c_v\,?
\]
Equivalently, the question asks whether such constants exist with
\[
 \forall B\in\mathcal B(f)\ \exists a_B\in\overline{\Q}^{\times},\qquad
 \langle f_B,f_B\rangle=a_B\prod_{v\in S(B)}c_v.
\]
An empty product is \(1\). The constants are allowed to lie in \(\C\), exactly as supplied; no additional nonvanishing hypothesis is asserted.

\emph{Unresolved qualifications.} The source leaves the pairing \(\langle\ ,\ \rangle\), its measure normalization, and the algebraic normalization of \(f_B\) unspecified. These choices matter: unrestricted complex rescaling of a representative can change its self-pairing by a nonalgebraic factor. The weight, level, character, and precise transfer hypotheses are also unspecified. Consequently, this is a conditional transcription of the intended period relation, requiring these data before it becomes a fully determined conjecture. The supplied dated assessment mentions a result under a technical discrete-series hypothesis and unresolved integral refinements, but supplies neither the hypothesis nor a formulation of those refinements.
\subsection{Short English statement}
For a fixed Hilbert newform over a totally real field, the question attributed to Shimura asks whether the self-pairings of its quaternionic Jacquet--Langlands transfers factor, up to nonzero algebraic multiples, into constants attached to the real places where the quaternion algebra splits. The admissible transfers and their normalizations require clarification. The supplied partially solved label is dated curation, not an independently verified conclusion.
\subsection{Sources}
\begin{itemize}
\item Supplied local UnsolvedMath registry record, permanent ID \(30001305\), source problem code \texttt{catalogue entry}, title: \emph{Shimura Period Relations for Jacquet--Langlands Transfers}. The record attributes the conjecture to Shimura, lists proposed year \(2009\), and supplies no named proposer.
\item Supplied report citation: \emph{Explicit Methods in Number Theory (2010)}, Oberwolfach Reports, pp.~1843--1920, DOI \texttt{10.4171/owr/2009/33}. \url{https://doi.org/10.4171/owr/2009/33}. The supplied background describes Oberwolfach Reports as documenting workshop discussions and open problems.
\item Supplied literature source labelled ``Prasanna Oberwolfach report'': \url{https://ems.press/content/serial-article-files/46233?nt=1}.
\item Supplied literature source labelled ``Quaternionic period results'': \url{https://arxiv.org/abs/1610.00166}.
\item Literature assessment dated exactly \texttt{2026-08-21}: supplied status ``partially solved.'' The supplied assessment says that Shimura's period relation is proved under a technical discrete-series hypothesis, while the extracted general formulation and integral refinements are not fully settled. This is supplied dated curation derived from the OWR record, not independently verified source consultation or a new literature assessment; its status and cited claims require independent verification.
\item The registry labels the statement \texttt{corrected\_verified} and describes its verification as meaning-preserving semantic, typographic, and LaTeX cleanup without unsupported source-dependent additions. This records supplied editorial provenance and does not establish mathematical correctness.
\end{itemize}
\end{document}
